% sections/introduction.tex — numbers only through \res{}; every claim has a row in paper/CLAIMS.md.
% Draft for the author's review (plan of sections approved 2026-10-18, PLAN §6: motivation, four contributions).
\section{Introduction}\label{sec:intro}

The mean--variance model of Markowitz~\cite{markowitz1952portfolio} with a limit on the number of assets held is a
mixed-integer quadratic programme, and the cardinality constraint makes it NP-hard~\cite{gao2013optimal}. Exact
methods based on perspective reformulations and outer approximation solve instances with hundreds of assets, mostly
with commercial solvers~\cite{mencarelli2018complex,frangioni2009computational,bertsimas2022scalable}, and heuristics
for the problem have a long tradition~\cite{chang2000heuristics,schaerf2002local}. Over the last decade the
cardinality-constrained mean--variance (CCMV) problem has also become a standard test case for formulations as
quadratic unconstrained binary optimisation (QUBO)~\cite{lucas2014ising,glover2019quantum} and for the samplers that
minimise them, from quantum and hybrid annealers to simulated annealing and tabu
search~\cite{herman2023quantum,venturelli2019reverse,mugel2022dynamic,lozano2026where}. A typical pipeline samples
subsets of assets from a binary surrogate of the problem and computes the weights of the best subsets by convex
optimisation~\cite{morapakula2025endtoend}. With simulated annealing or tabu search as the sampler, such a pipeline is
a classical heuristic in its own right and a standard baseline for quantum samplers.

Whether such pipelines are worth running is a question of comparison, and the published comparisons are hard to read.
First, positive results rest mostly on closed hybrid services and on weak, untuned or differently timed classical
baselines~\cite{mugel2022dynamic,lang2022strategic,kim2025quantum}, while studies that include strong classical
solvers find that these solve the portfolio problems considered~\cite{lozano2026where,stopfer2025quantum,chen2025benchmarking}.
Second, in the binary studies we reviewed the reference optimum is the optimum of the surrogate itself, so the loss
that the surrogate causes on the actual problem, with continuous weights, has not been
measured~\cite{brandhofer2022benchmarking,lozano2026where,lozano2026penalty}. Third, a local search on the true
objective with continuous weights~\cite{schaerf2002local,digaspero2011hybrid}, the classical counterpart of a pipeline that searches over
subsets, is absent from the comparisons we reviewed. Fourth, whether the cardinality constraint binds is not checked in
the closest recent benchmarks~\cite{mancilla2026constrained,lozano2026where,lozano2026penalty}, although discretised
portfolio problems with convex costs can be easy near the solution of their continuous
relaxation~\cite{rubiogarcia2022portfolio}, and on such instances a comparison of methods says little. Standards for
such benchmarks ask for the best classical algorithms, equal stopping criteria and budgets, and equal tuning
effort~\cite{barr1995designing,ronnow2014defining,dunning2018what,bucher2024towards,koch2026quantum}.
Finally, it is not clear that the accuracy of the optimiser matters for the portfolio out of sample, where the error
of the estimated inputs may dominate~\cite{michaud1989markowitz,gilli2011optimal}.

We address these points in a pre-registered benchmark with three research questions. RQ1 asks how far the
equal-weight binary surrogate limits a QUBO pipeline and how much the sampler adds. RQ2 asks whether, on instances where
the cardinality constraint binds and a cheap start is not already near-optimal, a tuned QUBO pipeline approaches the
optimum faster than open-source exact solvers and than a local search on the true objective. RQ3 asks whether the
accuracy of the optimiser matters for the portfolio out of sample. The paper makes four contributions.
\begin{enumerate}
\item A pre-registered benchmark~\cite{nosek2018preregistration}: the hypotheses, instances, metrics, tests and the
wording of every outcome were fixed and publicly time-stamped before the pilot runs, every later change is a dated
entry of a deviation log (\ref{app:deviations}), and every reported number is generated by code from registered runs.
\item An attribution of the error of the QUBO pipeline to the surrogate and to the sampler, measured against certified
optima and including the diversity of the pools that real samplers return, and a screening of trivial instances that
measures how often a combinatorial search is needed at all.
\item An anytime comparison of three classes of methods from a common start, with equal budgets and one tuning
protocol: QUBO pipelines with penalty and cardinality-preserving samplers, open-source exact solvers (SCIP with the
perspective reformulation and a multi-tree outer approximation with an open-source master problem), and a local
search on the true objective; together with a validation of the exact stack against published
times~\cite{bertsimas2022scalable}.
\item A test of whether optimisation accuracy matters out of sample, with alternative supports enumerated exactly by
cuts that exclude the supersets of better supports, their measured in-sample gaps and pre-registered equivalence tests.
\end{enumerate}

The main findings are as follows. At $N \le 24$ (RQ1), the ceiling of the surrogate rises with the dispersion of asset
volatilities and betas, and an oracle pool of the best subsets of the surrogate recovers most of the loss; the real
samplers, however, return few distinct subsets per call, and their pipelines stay above the target on
\res{e1.pipe.tabu.m200.nabove}--\res{e1.pipe.sa.m200.nabove} of \res{e1.pipe.ninstances} instances, against
\res{e1.eform.m50.all.nexceed} for the oracle pool. A cheap common start --- the continuous relaxation rounded to its
$K$ largest weights and refined --- is already within the target on
\res{e2.a3.factor.min.pct}--\res{e2.a3.factor.max.pct}\% of the synthetic factor instances, \res{e2.a3.c4.pct}\% of the
S\&P~500 bootstrap instances and \res{e2.a3.c5.pct}\% of the binding OR-Library points, and on the factor instances the
gaps that remain are at most \res{e2.a3.c2.gstart.max.bp}~bp of certainty equivalent per year. On the remaining
instances, the QUBO pipeline with a cardinality-preserving sampler approaches the optimum faster than the open-source
exact solvers (pooled ratio of primal integrals $R = \res{e2.h2}$), whereas the pipelines with the tuned penalty
samplers do not ($R = \res{e2.sa.h2}$ and \res{e2.tabu.h2}, post hoc). Against local search the pipeline is faster on
the typical factor instance but fails to reach the target on a minority of instances, and on average local search is
faster ($R_b = \res{e2.h2b}$; the size relative to the band of practical equivalence rests on one cell with three
clusters). The open-source exact representative needed in the median \res{e2.m11.t3.scip.median} and
\res{e2.m11.t4.scip.median} times the best published times on two validation sets, so the advantage of the pipeline over
exact methods is an advantage over open-source exact methods; it held when SCIP used aggressive primal heuristics and
when outer approximation represented the exact class in every cell (RQ2). Out of sample, the best non-optimal support, a few
basis points worse in sample, is equivalent to the optimum for minimum-variance portfolios, and the comparison is
inconclusive for mean--variance portfolios (RQ3).

This paper revisits our earlier study of the same pipeline~\cite{avdoshin2026discrete},
which compared open-source mixed-integer solvers, called through a modelling layer, with simulated annealing on the
binary surrogate. There, the gap of the pipeline was measured against the best solution found on a synthetic family
whose volatilities grow with the number of assets, the candidate sources of the gap were named but not separated, and
in the walk-forward the heuristic and the exact solver optimised different objectives. \ref{app:prior} lists how the
present design addresses these points; no result of the earlier study is reused, and its findings enter only as
exploratory hypotheses.

Section~\ref{sec:related} reviews related work, Section~\ref{sec:methods} defines the problem and the methods, and
Section~\ref{sec:design} the experimental design. Section~\ref{sec:results} reports the results,
Section~\ref{sec:discussion} discusses them, and Section~\ref{sec:conclusion} concludes.
